% Deep Learning -- Milestone 2 addendum (max 2 pages)
% Self-contained: compile separately from the main report.
% Uses the same CVPR style; intentionally avoids siunitx so it builds anywhere.
\documentclass[10pt,twocolumn,letterpaper]{article}
\usepackage{cvpr}
\usepackage{graphicx}
\usepackage{booktabs}
\usepackage{array}
\definecolor{cvprblue}{rgb}{0.21,0.49,0.74}
\usepackage[pagebackref,breaklinks,colorlinks,allcolors=cvprblue]{hyperref}
\newcommand{\TODO}[1]{\textbf{\color{red}[TODO: #1]}}

\def\paperID{*****}
\def\confName{CVPR}
\def\confYear{2026}

\title{Music Recommendation with Knowledge Graphs and GNNs\\
\large Deep Learning --- Milestone 2 Addendum}
\author{Group 7: Pedro Fanica (54346) \and
Aarthi Perumpillichira (66404) \and Maja Skwaro\'n (66083)}

\begin{document}
\maketitle

This addendum extends our Milestone-1 report with the training methods,
experiments and results produced since the previous submission. The data,
knowledge-graph construction and overall problem framing are unchanged; we
refer to the Milestone-1 report for those details.

\section{Training Methods}

We train three representation/recommendation components on top of the music
knowledge graph, and compare them against three baselines on a shared,
stratified user-level interaction split.

\paragraph{Symbolic-feature autoencoder.}
The 1{,}526-dimensional jSymbolic descriptor vector per track is standardised
and compressed by a fully connected autoencoder
($1526\!\rightarrow\!512\!\rightarrow\!256\!\rightarrow\!128$ with a symmetric
decoder), trained for 30 epochs with Adam (lr~$10^{-3}$, batch~256) to minimise
reconstruction MSE. The 128-d bottleneck is used as a dense \emph{content}
embedding for each track.

\paragraph{Knowledge-graph embeddings.}
We train RotatE and ComplEx (PyKEEN, 128-d complex / 256-d real, 100 epochs,
Adam) on the enriched graph with listening events excluded from the training
corpus, giving a structure-only entity representation.

\paragraph{Heterogeneous Graph Transformer (main model).}
Our principal recommender is an HGT over the typed graph (hidden~128, 4 heads,
2 layers, dropout~0.1; 100 epochs, Adam lr~$10^{-3}$, weight decay~$10^{-4}$),
with track nodes initialised from the autoencoder embeddings. It is trained
with a \emph{popularity-debiased listwise} loss (log-popularity prior,
$\lambda=0.2$, temperature~$0.1$) and scores items by user--track inner product
with seen items masked.

\paragraph{Baselines.}
MostPopular; user--user KNN-CF (cosine, $k$ chosen by validation sweep); and a
two-stage XGBoost hybrid (autoencoder-profile candidate generation followed by
\texttt{rank:ndcg} re-ranking on $[\text{user}\,\Vert\,\text{track}]$ content
features).

\section{Experiments}

All models share one user-level stratified train/val/test split (stratified by
genre and release decade); only interactions are partitioned, so song features
and graph structure are visible to all splits and no leakage occurs. We report
Recall@$K$, NDCG@$K$, Hit-Rate@$K$, MRR, catalogue Coverage and
Popularity-Bias@$K$ for $K\in\{5,10,20,50,100\}$, plus a composite Overall
score $0.6\,\mathrm{NDCG}+0.2\,\mathrm{Coverage}+0.2(1-\mathrm{PopBias})$.
Pairwise Wilcoxon and an omnibus Friedman test assess significance.
\TODO{State the train/val/test interaction counts.}

\section{Results}

\cref{tab:m2} reports top-10 test performance. Popularity is accurate but
nearly zero-coverage; KNN-CF wins the composite via far higher coverage and
lower popularity bias. \TODO{Fill the hybrid and HGT rows once the autoencoder
embeddings are regenerated, and state whether the GNN improves the
accuracy/diversity trade-off.}

\begin{table}[t]
\centering\footnotesize
\caption{Top-10 test performance (Cov.\ $=$ coverage).}
\label{tab:m2}
\begin{tabular}{lccccc}
\toprule
\textbf{Model} & Recall & NDCG & HitRate & Cov. & Overall \\
\midrule
MostPopular   & 0.1285 & 0.1143 & 0.2002 & 0.0064 & 0.1477 \\
KNN-CF        & 0.0923 & 0.0971 & 0.1540 & 0.9176 & 0.4185 \\
XGB-Hybrid    & \TODO{} & & & & \\
HGT (ours)    & \TODO{} & & & & \\
\bottomrule
\end{tabular}
\end{table}

\paragraph{Learning curves.}
\cref{fig:m2curves} shows the training diagnostics: autoencoder reconstruction
MSE (train/val), KGE loss, the HGT debiased-listwise loss with validation
Recall/NDCG, and the XGBoost per-round \texttt{ndcg@10}. \TODO{Insert the PNGs
exported to \texttt{data/final/}.}

\begin{figure}[t]
\centering
\fbox{\parbox[c][3.4cm][c]{0.92\linewidth}{\centering
  \TODO{Insert learning curves (AE / KGE / HGT / XGBoost).}}}
\caption{Learning curves for the trained models.}
\label{fig:m2curves}
\end{figure}

\paragraph{Significance \& qualitative.}
Accuracy gaps are statistically significant ($p<0.05$) but small in effect size
($d_z\!\approx\!0.03$--$0.11$); the decisive differences are in coverage and
popularity bias. Population- and user-level qualitative analyses
(genre/tempo/decade distribution shift, per-user case studies) are summarised
in the main report. \TODO{Add one qualitative figure if space permits.}

\section{Contribution of Each Member}
\TODO{Confirm the split below.}
\begin{itemize}\itemsep -2pt
  \item \textbf{Pedro Fanica:} \TODO{autoencoder, HGT + debiased loss,
  evaluation/qualitative framework, integration.}
  \item \textbf{Aarthi Perumpillichira:} \TODO{data linking, jSymbolic
  features, KNN/XGBoost baselines.}
  \item \textbf{Maja Skwaro\'n:} \TODO{KG construction, Wikidata/OWL
  enrichment, KGE baselines.}
\end{itemize}

\end{document}
